% Standalone problem document, generated from the adjacent README.md.
% Compile with XeLaTeX. Mathematical content is maintained in Markdown.
\documentclass[11pt,a4paper]{article}
\usepackage[fontsize=10.5pt]{scrextend}
\usepackage[margin=22mm,headheight=15pt,headsep=9mm,footskip=11mm]{geometry}
\usepackage{fontspec}
\setmainfont{texgyrepagella}[Extension=.otf,UprightFont=*-regular,BoldFont=*-bold,ItalicFont=*-italic,BoldItalicFont=*-bolditalic]
\setsansfont{texgyreheros}[Extension=.otf,UprightFont=*-regular,BoldFont=*-bold,ItalicFont=*-italic,BoldItalicFont=*-bolditalic]
\setmonofont{texgyrecursor}[Extension=.otf,UprightFont=*-regular,BoldFont=*-bold,ItalicFont=*-italic,BoldItalicFont=*-bolditalic,Scale=MatchLowercase]
\usepackage{amsmath,amssymb}
\usepackage{unicode-math}
\setmathfont{texgyrepagella-math.otf}
\usepackage{xcolor,microtype,enumitem,needspace,fancyhdr}
\usepackage{bookmark}
\definecolor{ink}{HTML}{17344B}
\definecolor{accent}{HTML}{197D80}
\definecolor{muted}{HTML}{596773}
\hypersetup{colorlinks=true,linkcolor=accent,urlcolor=accent,pdftitle={TR-12 - Classification
of perfect tensor formats with generic unique
decomposition},pdfauthor={Open Problems in Numerical Linear Algebra}}
\urlstyle{same}
\setlength{\parindent}{0pt}
\setlength{\parskip}{5pt plus 1pt minus 1pt}
\linespread{1.04}
\setlength{\emergencystretch}{3em}
\widowpenalty=10000
\clubpenalty=10000
\displaywidowpenalty=10000
\allowdisplaybreaks[1]
\setlist{nosep,leftmargin=1.5em}
\providecommand{\tightlist}{\setlength{\itemsep}{0pt}\setlength{\parskip}{0pt}}
\providecommand{\pandocbounded}[1]{#1}
\pagestyle{fancy}
\fancyhf{}
\fancyhead[L]{\sffamily\footnotesize\color{muted}OPEN PROBLEMS IN NUMERICAL LINEAR ALGEBRA}
\fancyhead[R]{\sffamily\bfseries\footnotesize\color{accent}TR-12}
\fancyfoot[L]{\parbox[b]{0.9\textwidth}{\sffamily\fontsize{8}{10}\selectfont\color{muted}Editorial ratings; a bounded search cannot certify that a problem remains unsolved.\\Literature check: 2026-09-08}}
\fancyfoot[R]{\sffamily\footnotesize\color{muted}\thepage}
\renewcommand{\headrulewidth}{0.3pt}
\renewcommand{\headrule}{\hbox to\headwidth{\color{accent}\leaders\hrule height \headrulewidth\hfill}}
\makeatletter
\renewcommand{\section}{\Needspace{5\baselineskip}\@startsection{section}{1}{0pt}{12pt plus 2pt minus 2pt}{4pt}{\sffamily\bfseries\large\color{ink}}}
\renewcommand{\subsection}{\Needspace{4\baselineskip}\@startsection{subsection}{2}{0pt}{9pt plus 1pt minus 1pt}{3pt}{\sffamily\bfseries\normalsize\color{ink}}}
\makeatother
\setcounter{secnumdepth}{0}
\begin{document}
{\sffamily\small\color{muted}Tensor computations\par}
\vspace{3pt}
{\raggedright\hyphenpenalty=10000\exhyphenpenalty=10000\sffamily\bfseries\fontsize{21}{24}\selectfont\color{ink}Classification
of perfect tensor formats with generic unique decomposition\par}
\vspace{7pt}
{\sffamily\small\textbf{Difficulty:} extreme\quad\textbf{Importance:} interesting
to the community\par}
\vspace{4pt}
{\color{accent}\hrule height 0.8pt}
\vspace{6pt}
\subsection{Statement}\label{statement}

Let \(d\ge3\), \(n_1,\ldots,n_d\ge2\), and suppose the dimension-count
value \[
r_*={\prod_{i=1}^d n_i\over 1+\sum_{i=1}^d(n_i-1)}
\] is an integer. Such a format is called perfect. The conjecture
asserts that a generic tensor in \(\bigotimes_{i=1}^d\mathbb C^{n_i}\)
has a unique minimal-length decomposition as a sum of rank-one tensors
if and only if, after permuting factors, the format is \[
(2,k,k)\ (k\ge2),\qquad(3,4,5),\qquad(2,2,2,3).
\] Generic means on a nonempty Zariski-open subset of the entire tensor
space. Rank one means \(v_1\otimes\cdots\otimes v_d\) with nonzero
complex factors. Two decompositions are the same when their rank-one
tensor summands agree after permutation; reciprocal rescalings within a
summand create no new decomposition. No symmetry constraints are
imposed.

\subsection{Relevance}\label{relevance}

This determines whether a decomposition algorithm can uniquely recover
components of a generic tensor at the parameter-count threshold. The
exceptional positive cases support explicit decomposition algorithms.

\subsection{References}\label{references}

\begin{enumerate}
\def\labelenumi{\arabic{enumi}.}
\tightlist
\item
  J. D. Hauenstein, L. Oeding, G. Ottaviani, and A. J. Sommese,
  \emph{Homotopy techniques for tensor decomposition and perfect
  identifiability}, J. reine angew. Math. 753 (2019), 1--22.
  \href{https://doi.org/10.1515/crelle-2016-0067}{DOI};
  \href{https://arxiv.org/pdf/1501.00090}{primary preprint}, §1,
  equation (2) and Conjecture 1.3, p.3; Theorems 1.1--1.2 establish the
  two sporadic positive cases.
\item
  A. Massarenti, M. Mella, and G. Staglianò, \emph{Effective
  identifiability criteria for tensors and polynomials}, J. Symbolic
  Comput. 87 (2018), 227--237.
  \href{https://www.iris.unict.it/retrieve/dfe4d22d-829e-bb0a-e053-d805fe0a78d9/Massarenti__Mella__Staglian\%C3\%B2_-_Effective_identi\%EF\%AC\%81ability_criteria_for_tensors_and_polynomials.pdf}{Author
  PDF}, §1, discussion of the Hauenstein et al.~generic-identifiability
  conjecture.
\end{enumerate}

\subsection{Status check --- 2026-09-08}\label{status-check-2026-09-08}

The latest originating preprint revision (November 9, 2016; journal
publication 2019) retains Conjecture 1.3, while replacing its former
symmetric analogue by Galuppi--Mella's theorem. Searches
\textit{"perfect identifiability" conjecture 2025 2026},
\textit{"perfect formats" "conjecture" "2025" OR "2026"}, and
\textit{"Hauenstein" "identifiability" "conjecture" "2026"} located
no resolution of this nonsymmetric statement. This is a bounded search
without a new 2026 reaffirmation.
\end{document}
